\begin{abstract}
Este trabajo propone un método para detectar esquemas de \emph{pump \& dump}
en mercados de criptoactivos usando teoría espectral de grafos. La idea
central es modelar el mercado como una red donde los nodos son tokens y las
aristas representan correlaciones entre sus precios. Primero filtramos el
ruido de mercado con Teoría de Matrices Aleatorias (RMT), usando dos ventanas
de tiempo: una larga para calibrar los parámetros del filtro y una corta para
detectar anomalías. Luego reducimos la red para que no sea demasiado densa,
usando un árbol de expansión mínima (MST). Sobre la red resultante calculamos
el \emph{Inverse Participation Ratio} (IPR), que mide qué tan concentrada está
la energía de cada modo de vibración de la red. La hipótesis es que cuando un
grupo de tokens es manipulado coordinadamente, forma una comunidad artificial
que se detecta como un aumento anómalo del IPR. Para evitar falsos positivos,
comparamos el IPR observado contra el que produciría una red aleatoria con la
misma estructura de conexiones (\emph{Configuration Model}). La validación se
realiza sobre ${\approx}351$ eventos confirmados del dataset de La~Morgia
\textit{et al.}~\cite{lamorgia2020}, usando datos de 1~minuto de la API pública
de Binance para pares \texttt{SYM/BTC} durante 2021--2022.
\end{abstract}

\begin{IEEEkeywords}
Teoría de grafos, análisis espectral, Laplaciana normalizada, Teoría de
Matrices Aleatorias, distancia de Mantegna, esparsificación MST, \emph{Configuration
Model}, criptoactivos, detección de manipulación de mercado, \emph{pump and
dump}, \emph{Inverse Participation Ratio}.
\end{IEEEkeywords}
